% !TeX root = ../../book.tex
\begin{chapex}
视频共享网站 \textit{YouTube} 为每个视频分配一个唯一标识符，该标识符是由 11 个字符组成的字符串，取自集合
\[ \{ \mathtt{A}, \mathtt{B}, \dots, \mathtt{Z}, \mathtt{a}, \mathtt{b}, \dots, \mathtt{z}, \mathtt{0}, \mathtt{1}, \mathtt{2}, \mathtt{3}, \mathtt{4}, \mathtt{5}, \mathtt{6}, \mathtt{7}, \mathtt{8}, \mathtt{9}, \text{\texttt{-}}, \text{\texttt{\_}} \} \]
这个字符串实际上是一个用 base-64 表示的自然数，其中上面集合中的字符按顺序代表数字$0$到$63$---$\mathtt{C}$代表$2$, $\ mathtt{c}$ 代表 $28$, $\mathtt{3}$ 代表 $55$, 而 \texttt{\_} 代表 $63$。根据这个规则，求 base-64 位展开为 $\mathtt{dQw4w9WgXcQ}$ 的自然数，求自然数 $7\,159\,047\,702\,620\,056\,984$ 的 base-64 展开。
\end{chapex}

\begin{chapex}
设 $a, b, c, d \in \mathbb{Z}$。什么条件下 $(a+b\sqrt{2})(c+d\sqrt{2})$ 是一个整数？
\end{chapex}

\begin{chapex}
假设整数 $m$ 除以 $3$ 余数为 $i$，整数 $n$ 除以 $3$ 余数为 $j$，其中 $0 \le i,j < 3$。证明 $m+n$ 除以 $3$ 与 $i+j$ 除以 $3$ 余数相等。
\end{chapex}

\begin{chapex}
$n^2$ 除以 $3$ 可能的余数是多少？其中 $n \in \mathbb{Z}$。
\end{chapex}

\begin{definition}
如果只要 $a$ 和 $b$ 是集合 $X$ 的元素，$a \odot b$ 也是集合 $X$ 的元素，则集合 $X$ 在运算 $\odot$ 下是\textbf{封闭的}。
\end{definition}

在 \Crefrange{cqClosureOfNumberSetsBegin}{cqClosureOfNumberSetsEnd} 中，确定数集 $\mathbb{N}$, $\mathbb{Z}$, $\mathbb{Q}$ 和 $\mathbb{R}$ 哪个在问题定义的 $\odot$ 运算下是封闭的。

\begin{multicols}{2}
\begin{chapex}
\label{cqClosureOfNumberSetsBegin}
$a \odot b = a + b$
\end{chapex}

\begin{chapex}
$a \odot b = a - b$
\end{chapex}

\begin{chapex}
$a \odot b = (a-b)(a+b)$
\end{chapex}

\begin{chapex}
$a \odot b = (a-1)(b-1) + 2(a+b)$
\end{chapex}

\begin{chapex}
$a \odot b = \dfrac{a}{b^2+1}$
\end{chapex}

\begin{chapex}
$a \odot b = \dfrac{a}{\sqrt{b^2+1}}$
\end{chapex}

\begin{chapex}
\label{cqClosureOfNumberSetsEnd}
$a \odot b = \begin{cases} a^{b} & \text{if $b > 0$} \\ 0 & \text{if $b \not\in \mathbb{Q}$} \end{cases}$
\end{chapex}
\end{multicols}

\begin{definition}
如果对于 $\mathbb{Q}$ 上的某个非零多项式 $p(x)$, $p(\alpha) = 0$，则复数 $\alpha$ 是\textbf{代数的}。
\end{definition}

\begin{chapex}
设 $x$ 是有理数。证明 $x$ 是代数数。
\end{chapex}

\begin{chapex}
证明 $\sqrt{2}$ 是代数数。
\end{chapex}

\begin{chapex}
证明 $\sqrt{2} + \sqrt{3}$ 是代数数。
\end{chapex}

\begin{chapex}
$x$ 和 $y$ 是任意两个有理数，证明 $x+yi$ 是代数数。
\end{chapex}

\subsection*{判断题}

\tfquestiontext{cqGettingStartedTFBegin}{cqGettingStartedTFEnd}

\begin{chapex} % False
\label{cqGettingStartedTFBegin}
所有整数都是自然数。
\end{chapex}

\begin{chapex} % True
所有整数都是有理数。
\end{chapex}

\begin{chapex} % True
所有整数都能整除0。
\end{chapex}

\begin{chapex} % True
所有整数都能整除它的平方。
\end{chapex}

\begin{chapex} % True
有理数的平方是有理数。
\end{chapex}

\begin{chapex} % False
正有理数的平方根是有理数。
\end{chapex}

\begin{chapex} % False
当整数 $a$ 除以正整数 $b$，余数总是小于 $a$。
\end{chapex}

\begin{chapex} % False
\label{cqGettingStartedTFEnd}
二次多项式都有两个不同的复数根。
\end{chapex}

\subsection*{可能性问题}

\asnquestiontext{cqGettingStartedASNBegin}{cqGettingStartedASNEnd}

\begin{chapex} % Always
\label{cqGettingStartedASNBegin}
设 $n,b_1,b_2 \in \mathbb{N}$ 其中 $1 < n < b_1 < b_2$。则 $n$ 的 base-$b_1$ 展开等于 $n$ 的 base-$b_2$ 展开。
\end{chapex}

\begin{chapex} % Never
设 $n,b_1,b_2 \in \mathbb{N}$ 其中 $1 < b_1 < b_2 < n$。则 $n$ 的 base-$b_1$ 展开等于 $n$ 的 base-$b_2$ 展开。
\end{chapex}

\begin{chapex} % Sometimes
设 $a,b,c \in \mathbb{Z}$ 并假设 $a$ 整除 $c$ 且 $b$ 整除 $c$。则 $ab$ 整除 $c$。
\end{chapex}

\begin{chapex} % Always
设 $a,b,c \in \mathbb{Z}$ 并假设 $a$ 整除 $c$ 且 $b$ 整除 $c$。则 $ab$ 整除 $c^2$。
\end{chapex}

\begin{chapex} % Always
设 $x,y \in \mathbb{Q}$ 并设 $a,b,c,d \in \mathbb{Z}$ 其中 $cy+d \ne 0$。则 $\dfrac{ax+b}{cy+d} \in \mathbb{Q}$。
\end{chapex}

\begin{chapex} % Sometimes
设 $\dfrac{a}{b}$ 是一个有理数。则 $a \in \mathbb{Z}$ 且 $b \in \mathbb{Z}$。
\end{chapex}

\begin{chapex} % Sometimes
设 $x \in \mathbb{R}$ 并假设 $x^2 \in \mathbb{Q}$。则 $x \in \mathbb{Q}$。
\end{chapex}

\begin{chapex} % Always
设 $x \in \mathbb{R}$ 并假设 $x^2+1 \in \mathbb{Q}$ 且 $x^5+1 \in \mathbb{Q}$。 则 $x \in \mathbb{Q}$。
\end{chapex}

\begin{chapex} % Always
\label{cqGettingStartedASNEnd}
设 $p(x) = ax^2+bx+c$ 是一个多项式， 其中 $a,b,c \in \mathbb{R}$ 且 $a \ne 0$，假设 $u+vi$ 是 $p(x)$ 的一个复数根,其中 $v \ne 0$。则 $u-vi$ 是 $p(x)$ 的根。
\end{chapex}